\PurePy is a subset of Python: every well-formed program is a syntactically valid Python program, and if \PurePy
gives a run a result, that run produces the same result under Python. \paperonly{\figref{taxonomy} shows the
relationship between the two languages, with the guarantee appearing as containment.}

Well-formedness is a property of the program, decided by \PurePy's static checks. Semantic validity is
a property of a run rather than a program, because some of the Python behaviours \PurePy excludes cannot be rejected
statically (\secref{no-coercions}). An implementation may report a dynamically excluded operation by raising
\pyinline{UnsupportedOperationError}, but is not obliged to. The guarantee covers only valid runs, and a
stock Python interpreter will simply produce the Python result. \added{An implementation conforms if it accepts
exactly the well-formed programs and, on every valid run, produces the Python result.} The conformance suite is
organised the same way\paperonly{ (\secref{reference-impl})}\speconly{ (\secref{testing})}.
